% chromatic-poly.tex

%%%%%%%%%%%%%%%
\begin{frame}{}
	\begin{center}
		\red{$P(G, k)$: How many ways to color the vertices of $G$ with $k$ colors?}

	  % \fig{width = 0.70\textwidth}{figs/3-vertex-coloring}

		\pause
		\[
			P(P_{4}, k) = \pause k(k-1)^3
		\]
		\pause
		\[
			P(C_4, k) = \pause k(k-1)^2(k-2)
		\]
		\pause
		\[
			P(K_4, k) = \pause k(k-1)(k-2)(k-3)
		\]
	\end{center}
\end{frame}
%%%%%%%%%%%%%%%

%%%%%%%%%%%%%%%
\begin{frame}{}
	\begin{theorem}[Deletion-Contraction Recurrence]
		Let $G$ be a simple graph and $e$ be an edge of $G$. Then
		\[
			P(G, k) = \purple{P(G - e, k)} - \violet{P(G / e, k)},
		\]
		where \purple{$G - e$} is the graph obtained from $G$ by \purple{deleting} $e$, \\
		and \violet{$G / e$} is the graph obtained from $G$ by \violet{contracting} $e$.
	\end{theorem}

	\pause
	\fig{width = 0.70\textwidth}{figs/chromatic-poly-example}
	\pause
	\[
		k(k-1)(k-2)(k-3) = k(k-1)(k-2)^{2} - k(k-1)(k-2)
	\]
\end{frame}
%%%%%%%%%%%%%%%

%%%%%%%%%%%%%%%
\begin{frame}{}
	\begin{theorem}[Deletion-Contraction Recurrence]
		\[
			P(G, k) = \purple{P(G - e, k)} - \violet{P(G / e, k)}
		\]
	\end{theorem}

	\fig{width = 0.70\textwidth}{figs/chromatic-poly-example}

	\pause
	\[
		P(G - e, k) =  P(G, k) + P(G / e, k)
	\]
	\pause
	\vspace{-0.30cm}
	\begin{align*}
		\purple{P(G - e, k)} =\; &P(G - e, k) \text{ with $v$ and $w$ having \red{\it different} colors} \\
		  +\; &P(G - e, k) \text{ with $v$ and $w$ having the \red{\it same} color}
	\end{align*}
\end{frame}
%%%%%%%%%%%%%%%

%%%%%%%%%%%%%%%
\begin{frame}{}
	\[
		P(G, k) = \purple{P(G - e, k)} - \violet{P(G / e, k)}
	\]

	\fig{width = 0.65\textwidth}{figs/chromatic-poly-example-2}

	\pause
	\vspace{-0.30cm}
	\begin{align*}
		P(G, k) &= k(k-1)^4 - 3k(k-1)^3 + 2k(k-1)^2 + k(k-1)(k-2) \\
		        &= k^5 - 7k^4 + 18k^3 - 20k^2 + 8k
	\end{align*}

	\pause
	\begin{theorem}[Chromatic Polynomial]
		$P(G, k)$ is a polynomial in $k$.
	\end{theorem}

	\pause
	\vspace{-0.30cm}
	\[
		\red{k^n} \blue{-} \brown{m} k^{n - 1} \blue{+}
		  \cdots \blue{-} \cdots + ak + \cyan{0}
	\]
\end{frame}
%%%%%%%%%%%%%%%